\chapter{Обучение с временной разностью}
\label{ch:td}

\begin{quote}
\itshape
Если и можно выделить одну идею как центральную и новаторскую в обучении с подкреплением, то это, несомненно, обучение с временной разностью (TD).
\end{quote}

\bigskip

Обучение с временной разностью --- наиважнейшая идея в обучении с подкреплением.  Оно синтезирует два мощных направления мысли: \emph{выборочный} подход методов Монте-Карло, который обучается непосредственно из опыта без модели среды, и \emph{бутстрэппинговый} подход динамического программирования, который обновляет оценки на основе других оценок, а не дожидаясь окончательного результата.  В итоге получается семейство алгоритмов, способных обучаться онлайн, после каждого отдельного перехода, как в эпизодических, так и в продолжающихся задачах, не требуя модели среды.

В этой главе семейство TD развивается во всей полноте.  Мы начинаем с простейшего его члена, TD(0), и устанавливаем его ключевые свойства: ошибку TD, свойство нулевого среднего при истинной функции ценности и сходимость.  Затем переходим к TD-алгоритмам управления\,---\,SARSA, Q-learning, Expected SARSA и Double Q-learning\,---\,и тщательно сравниваем их характер применительно к текущей и другой стратегии.  Обобщаем до $n$-шаговых отдач, которые интерполируют между TD и Монте-Карло, а затем до полного алгоритма TD($\lambda$) со следами приемлемости, который одновременно объединяет информацию со всех горизонтов.  Завершаем обобщённой оценкой преимущества (GAE) --- современным методом для оценки преимущества в методах градиента стратегии, который прямо основан на механизме TD($\lambda$), --- и тщательным анализом смещения, дисперсии и сходимости по всему семейству.

% ===========================================================
\section{Предсказание TD(0)}
\label{sec:td0}

\subsection{От метода Монте-Карло к временной разности}

Вспомним из главы~\ref{ch:mc}, что предсказание методом Монте-Карло (МК) оценивает функцию ценности $V^\pi$, усредняя полные выборочные отдачи.  После каждого полного эпизода каждое посещённое состояние $S_t$ получает обновление
\begin{equation}
\label{eq:mc-update}
V(S_t) \leftarrow V(S_t) + \alpha\bigl(G_t - V(S_t)\bigr),
\end{equation}
где $G_t = \sum_{k=0}^{\infty}\gamma^k R_{t+k+1}$ --- дисконтированная отдача с момента $t$.  Этот оценщик несмещён\,---\,$\E_\pi[G_t \mid S_t = s] = V^\pi(s)$\,---\,но имеет три ограничения: требует завершения эпизода до начала обучения; накапливает дисперсию вознаграждений на длинных горизонтах; вовсе не применим к продолжающимся (не эпизодическим) задачам.

Уравнение Беллмана\index{Bellman equation} для фиксированной стратегии $\pi$ предлагает альтернативную цель.  Поскольку
\begin{equation}
\label{eq:bellman-v-ch6}
V^\pi(s) = \E_\pi\bigl[R_{t+1} + \gamma V^\pi(S_{t+1}) \mid S_t = s\bigr],
\end{equation}
истинная ценность состояния $s$ равна ожидаемому немедленному вознаграждению плюс дисконтированная ценность следующего состояния.  Уравнение~\eqref{eq:bellman-v-ch6} точно при $V = V^\pi$, но даже при несовершенной оценке $V \approx V^\pi$ величина
\begin{equation}
\label{eq:td0-target}
R_{t+1} + \gamma V(S_{t+1})
\end{equation}
представляет собой полезную \emph{одношаговую цель}\index{one-step target}, которую можно вычислить сразу после наблюдения перехода $(S_t, A_t, R_{t+1}, S_{t+1})$.  Заменяя полную отдачу $G_t$ в~\eqref{eq:mc-update} этой одношаговой целью, получаем

\begin{equation}
\label{eq:td0-update}
\boxed{V(S_t) \leftarrow V(S_t) + \alpha\Bigl(R_{t+1} + \gamma V(S_{t+1}) - V(S_t)\Bigr).}
\end{equation}
Это алгоритм \textbf{TD(0)}\index{TD(0)}, также называемый одношаговым методом временной разности.

\subsection{Ошибка TD}

\begin{definition}[Ошибка TD]
\label{def:td-error}
\index{TD error}
\textbf{Ошибка временной разности} в момент $t$ равна
\begin{equation}
\label{eq:td-error}
\delta_t = R_{t+1} + \gamma V(S_{t+1}) - V(S_t).
\end{equation}
\end{definition}

Ошибка TD измеряет расхождение между текущей оценкой $V(S_t)$ и одношаговой целью $R_{t+1}+\gamma V(S_{t+1})$.  С этим обозначением обновление TD(0)~\eqref{eq:td0-update} упрощается до
\[
V(S_t) \leftarrow V(S_t) + \alpha\,\delta_t.
\]
Знак $\delta_t$ имеет непосредственную интуитивную интерпретацию:
\begin{itemize}[leftmargin=*]
  \item $\delta_t > 0$: наблюдаемый переход оказался лучше, чем ожидалось; $V(S_t)$ следует увеличить.
  \item $\delta_t < 0$: состояние было переоценено; $V(S_t)$ следует уменьшить.
  \item $\delta_t = 0$: текущая оценка локально согласована с уравнением Беллмана.
\end{itemize}

\begin{definition}[Бутстрэппинг]
\label{def:bootstrap}
\index{bootstrapping}
\textbf{Бутстрэппинг} --- это использование текущей оценки функции ценности как части цели в правиле обновления.  В TD(0) цель $R_{t+1}+\gamma V(S_{t+1})$ зависит от оценки $V$ в следующем состоянии, поэтому метод использует бутстрэппинг.  Монте-Карло не использует бутстрэппинг, поскольку $G_t$ является истинной отдачей, независимой от $V$.
\end{definition}

\subsection{Свойство нулевого среднего ошибки TD}

Следующий результат является фундаментальным: он показывает, что $\delta_t$ --- корректный (нулевого среднего) градиентный сигнал, когда функция ценности точна.

\begin{theorem}[Нулевое среднее ошибки TD]
\label{thm:td-error-zero-mean}
\index{TD error!zero-mean property}
Если $V = V^\pi$, то для каждого состояния $s$:
\begin{equation}
\label{eq:td-zero-mean}
\E_\pi[\delta_t \mid S_t = s] = 0.
\end{equation}
\end{theorem}

\begin{proof}
Подставляя $V = V^\pi$ в определение~\eqref{eq:td-error}:
\[
\delta_t = R_{t+1} + \gamma V^\pi(S_{t+1}) - V^\pi(S_t).
\]
Беря условное математическое ожидание при $S_t = s$:
\begin{align*}
\E_\pi[\delta_t \mid S_t = s]
&= \E_\pi\bigl[R_{t+1} + \gamma V^\pi(S_{t+1}) \mid S_t = s\bigr] - V^\pi(s)\\
&= V^\pi(s) - V^\pi(s) = 0,
\end{align*}
где вторая строка использует уравнение Беллмана~\eqref{eq:bellman-v-ch6}.
\end{proof}

Теорема~\ref{thm:td-error-zero-mean} подразумевает, что в неподвижной точке TD-обновление имеет нулевое ожидаемое изменение.  До сходимости $\delta_t \ne 0$ в среднем, и его знак последовательно указывает направление, уменьшающее ошибку Беллмана.

\subsection{Алгоритм TD(0)}

\begin{algorithm}[ht]
\caption{TD(0) для оценивания стратегии}
\label{alg:td0}
\KwIn{Стратегия $\pi$, шаг $\alpha\in(0,1]$, дисконт $\gamma\in[0,1)$}
Инициализировать $V(s) \leftarrow 0$ для всех $s \in \cS$, $V(s_{\mathrm{term}}) \leftarrow 0$\;
\For{каждый эпизод}{
  Инициализировать $S_0$\;
  \For{$t = 0, 1, 2, \ldots$ пока $S_t$ не терминальное}{
    Выбрать $A_t \sim \pi(\cdot \mid S_t)$\;
    Выполнить действие $A_t$; наблюдать $R_{t+1}$ и $S_{t+1}$\;
    $\delta_t \leftarrow R_{t+1} + \gamma V(S_{t+1}) - V(S_t)$\;
    $V(S_t) \leftarrow V(S_t) + \alpha\,\delta_t$\;
  }
}
\end{algorithm}

Заметим, что, в отличие от Монте-Карло, обновление в строке~8 происходит \emph{немедленно} после каждого перехода\,---\,ждать конца эпизода не нужно.

\subsection{Сходимость табличного TD(0)}

Сходимость TD(0) можно понять в рамках оператора Беллмана\index{Bellman operator}.  Для фиксированной стратегии $\pi$ определим
\[
(T^\pi V)(s) = \E_\pi\bigl[R_{t+1} + \gamma V(S_{t+1}) \mid S_t = s\bigr].
\]
Уравнение Беллмана принимает вид $V^\pi = T^\pi V^\pi$.  Поскольку $T^\pi$ является $\gamma$-сжатием в супремум-норме, его единственная неподвижная точка --- $V^\pi$.  TD(0) можно интерпретировать как стохастическое приближение итерации $V \leftarrow T^\pi V$: вместо вычисления полного математического ожидания используется один выборочный переход.

\begin{theorem}[Сходимость табличного TD(0)]
\label{thm:td0-convergence}
\index{TD(0)!convergence}
В табличном случае, если $\pi$ фиксирована, все состояния посещаются бесконечно часто, вознаграждения ограничены, и шаги $\{\alpha_t(s)\}$ удовлетворяют условиям Роббинса--Монро\index{Robbins--Monro conditions}:
\begin{equation}
\label{eq:rm-conditions}
\sum_{t=1}^{\infty}\alpha_t(s) = \infty, \qquad \sum_{t=1}^{\infty}\alpha_t^2(s) < \infty \quad \mbox{для каждого } s,
\end{equation}
то $V_t(s) \to V^\pi(s)$ почти наверное для всех $s\in\cS$.
\end{theorem}

Условия Роббинса--Монро~\eqref{eq:rm-conditions} требуют, чтобы шаги были достаточно велики для коррекции любой начальной ошибки (первое условие), но достаточно малы для усреднения шума (второе условие).  Расписание $\alpha_t(s) = 1/N_t(s)$, где $N_t(s)$ --- счётчик посещений, удовлетворяет обоим условиям.  Постоянный шаг $\alpha$ удовлетворяет только первому, поэтому он не сходится к точному $V^\pi$, но отслеживает его в нестационарной среде.

\subsection{Разобранный пример: случайное блуждание}

\begin{example}[Случайное блуждание по пяти состояниям]
\label{ex:random-walk}
\index{random walk example}
Рассмотрим цепь Маркова с пятью состояниями $A, B, C, D, E$, расположенными линейно, плюс два терминальных состояния $L$ (слева) и $R$ (справа).  Из любого нетерминального состояния агент движется влево или вправо с равной вероятностью $0.5$.  Единственное ненулевое вознаграждение --- $+1$ при переходе в правое терминальное состояние $R$; все остальные вознаграждения равны $0$.  Коэффициент дисконтирования $\gamma = 1$.  Истинные ценности состояний (то есть вероятность в конечном счёте достичь $R$) равны:
\[
V^\pi(A)=\tfrac{1}{6}, \quad V^\pi(B)=\tfrac{2}{6}, \quad V^\pi(C)=\tfrac{3}{6}, \quad V^\pi(D)=\tfrac{4}{6}, \quad V^\pi(E)=\tfrac{5}{6}.
\]
Предположим, что инициализируем $V(s) = 0.5$ для всех состояний и запускаем TD(0) с $\alpha = 0.1$.  Рассмотрим одну траекторию $C \to D \to E \to R$.

\textbf{Шаг 1}: $S_t = C$, $A_t = \mbox{вправо}$, $R_{t+1} = 0$, $S_{t+1} = D$.
\[
\delta_t = 0 + 1.0 \cdot 0.5 - 0.5 = 0.
\]
На этом шаге обновлений нет (поскольку $V(D) = V(C) = 0.5$ изначально).

\textbf{Шаг 2}: $S_t = D$, $R_{t+1} = 0$, $S_{t+1} = E$.
\[
\delta_t = 0 + 0.5 - 0.5 = 0.
\]
Снова без изменений.

\textbf{Шаг 3}: $S_t = E$, $R_{t+1} = 1$, $S_{t+1} = R$ (терминальное, $V(R)=0$).
\[
\delta_t = 1 + 0 - 0.5 = 0.5.
\]
\[
V(E) \leftarrow 0.5 + 0.1 \times 0.5 = 0.55.
\]
Сигнал вознаграждения распространяется только на один шаг назад.  На протяжении многих эпизодов TD(0) постепенно распространяет оценки ценности назад по цепочке.  Примерно после 100 эпизодов оценки близки к истинным значениям.  Этот пример демонстрирует как скоростное преимущество TD (обучение внутри каждого эпизода), так и проблему присвоения кредита: каждый новый сигнал вознаграждения должен распространяться назад по пространству состояний шаг за шагом.
\end{example}

% ===========================================================
\section{Преимущества методов TD}
\label{sec:td-advantages}

Методы TD занимают привилегированное положение в области обучения с подкреплением, поскольку унаследовали преимущества как от Монте-Карло, так и от динамического программирования, избегая их главных слабостей.  Перечислим эти преимущества.

\paragraph{Не требуют полного эпизода.}
Методы Монте-Карло должны ждать конца эпизода для вычисления $G_t$.  Это проблематично для длинных эпизодов и невозможно для продолжающихся задач без естественного завершения.  TD(0) обновляется после каждого отдельного шага, что делает его немедленно применимым к любой структуре задач.

\paragraph{Работают для продолжающихся задач.}
Поскольку цель TD $R_{t+1}+\gamma V(S_{t+1})$ всегда доступна (терминальное состояние не нужно), методы TD естественно обрабатывают продолжающиеся задачи.  Методы Монте-Карло в своей базовой форме не могут применяться в таких условиях.

\paragraph{Меньше дисперсия, чем у Монте-Карло.}
Отдача Монте-Карло $G_t = R_{t+1} + \gamma R_{t+2} + \gamma^2 R_{t+3} + \cdots$ накапливает дисперсию от каждого будущего вознаграждения.  Цель TD $R_{t+1}+\gamma V(S_{t+1})$ включает лишь одно случайное вознаграждение и (относительно гладкую) оценку $V(S_{t+1})$.  Формально, если каждое вознаграждение имеет дисперсию $\sigma^2$ и переходы независимы, дисперсия МК растёт примерно как $\sum_{k=0}^{T-t-1} \gamma^{2k}\sigma^2$, тогда как дисперсия TD примерно равна $\sigma^2 + \gamma^2 \Var[V(S_{t+1})]$.

\paragraph{Компромисс смещение--дисперсия.}\index{bias--variance trade-off}
Ценой меньшей дисперсии является то, что цель TD \emph{смещена}: поскольку в начале обучения $V \ne V^\pi$, цель $R_{t+1}+\gamma V(S_{t+1})$ не равна $V^\pi(S_t)$ в математическом ожидании.  По теореме~\ref{thm:td-error-zero-mean} смещение исчезает при сходимости.  Монте-Карло несмещён, но имеет высокую дисперсию.  Этот фундаментальный компромисс\,---\,бутстрэппинг покупает меньшую дисперсию ценой смещения\,---\,присутствует во всех TD-методах и количественно описан в разделе~\ref{sec:td-bias-variance}.

\paragraph{Более быстрое присвоение кредита.}
Поскольку TD распространяет оценки ценности на один шаг в каждом эпизоде, полезная информация из вознаграждаемых переходов включается в более ранние состояния быстрее, чем при МК.  $n$-шаговое TD (раздел~\ref{sec:nstep-td}) и следы приемлемости (раздел~\ref{sec:td-lambda}) предоставляют механизмы для эффективного распространения кредита на несколько шагов.

\paragraph{Пакетное TD(0) и оценка максимального правдоподобия.}\index{batch TD(0)}
Когда доступен фиксированный пакет эпизодов, можно многократно применять обновления TD(0) до сходимости.  Это \emph{пакетное TD(0)} сходится к \emph{оценке максимального правдоподобия} (оценке с определённостью эквивалентности) $V^\pi$ по данным\,---\,точному решению, которое вычислялось бы построением эмпирического МПР из пакета и его решением.  Для сравнения, пакетный Монте-Карло сходится к оценке, минимизирующей среднеквадратичную ошибку на тренировочных отдачах.  В задачах с ограниченными данными TD-решение часто лучше обобщается, поскольку использует марковскую структуру.

% ===========================================================
\section{SARSA: управление TD на текущей стратегии}
\label{sec:sarsa}

\subsection{От предсказания к управлению}

Оценивание стратегии говорит нам, насколько хороша текущая стратегия.  Управление требует её улучшения.  Как и в управлении методом Монте-Карло (глава~\ref{ch:mc}), ключевое наблюдение состоит в том, что нужно оценивать \emph{функцию ценности действия}\index{action-value function} $Q^\pi(s,a) = \E_\pi[G_t \mid S_t=s, A_t=a]$, а не $V^\pi(s)$, поскольку жадное улучшение требует сравнения действий:
\[
\pi'(s) \in \argmax_{a\in\cA}\; Q(s,a).
\]
Уравнение Беллмана для $Q^\pi$ имеет вид:
\begin{equation}
\label{eq:bellman-q-ch6}
Q^\pi(s,a) = \E_\pi\bigl[R_{t+1} + \gamma Q^\pi(S_{t+1},A_{t+1}) \mid S_t=s, A_t=a\bigr].
\end{equation}
Заменяя математическое ожидание в~\eqref{eq:bellman-q-ch6} одной выборкой и применяя полученное обновление, получаем:

\subsection{Обновление SARSA}

\begin{definition}[SARSA]
\label{def:sarsa}
\index{SARSA}
\textbf{SARSA}\footnote{Название происходит от пятёрки $(S_t, A_t, R_{t+1}, S_{t+1}, A_{t+1})$, используемой в каждом обновлении.} (Раммери и Нираньян, 1994; Саттон, 1996) --- это алгоритм управления TD на текущей стратегии с правилом обновления:
\begin{equation}
\label{eq:sarsa}
Q(S_t,A_t) \leftarrow Q(S_t,A_t) + \alpha\Bigl(R_{t+1} + \gamma Q(S_{t+1},A_{t+1}) - Q(S_t,A_t)\Bigr).
\end{equation}
\end{definition}

Критическая особенность SARSA состоит в том, что $A_{t+1}$ извлекается из \emph{той же} стратегии $\pi$, которая используется для сбора данных.  Это делает SARSA \textbf{алгоритмом на текущей стратегии}\index{on-policy}: он оценивает и улучшает стратегию, которая генерирует данные.

\begin{proposition}
Цель $Y_t^{\mathrm{SARSA}} = R_{t+1} + \gamma Q(S_{t+1},A_{t+1})$ является несмещённой выборкой оператора Беллмана, применённого к $Q$:
\[
\E_\pi\bigl[Y_t^{\mathrm{SARSA}} \mid S_t=s, A_t=a\bigr] = (T^\pi Q)(s,a),
\]
где $(T^\pi Q)(s,a) = \E_\pi[R_{t+1}+\gamma Q(S_{t+1},A_{t+1}) \mid S_t=s, A_t=a]$.
\end{proposition}

\begin{proof}
По определению условного математического ожидания утверждение следует непосредственно.
\end{proof}

\begin{algorithm}[ht]
\caption{SARSA (управление TD на текущей стратегии)}
\label{alg:sarsa}
\KwIn{Шаг $\alpha$, дисконт $\gamma$, исследование $\varepsilon$}
Инициализировать $Q(s,a)$ для всех $s,a$; $Q(s_{\mathrm{term}},\cdot)\leftarrow 0$\;
\For{каждый эпизод}{
  Инициализировать $S$\;
  Выбрать $A$ из $\pi(\cdot\mid S)$ на основе $Q$ (например, $\varepsilon$-жадная)\;
  \For{каждый шаг эпизода пока $S$ не терминальное}{
    Выполнить действие $A$; наблюдать $R, S'$\;
    Выбрать $A'$ из $\pi(\cdot\mid S')$ (та же $\varepsilon$-жадная)\;
    $Q(S,A) \leftarrow Q(S,A) + \alpha\bigl(R + \gamma Q(S',A') - Q(S,A)\bigr)$\;
    $S \leftarrow S'$;\quad $A \leftarrow A'$\;
  }
}
\end{algorithm}

\subsection{Сходимость SARSA}

\begin{theorem}[Сходимость SARSA]
\index{SARSA!convergence}
В табличном случае SARSA сходится к $Q^*$ почти наверное, если: (1) все пары состояние-действие посещаются бесконечно часто; (2) стратегия сходится к жадной (например, $\varepsilon_t \to 0$); (3) шаги удовлетворяют условиям Роббинса--Монро~\eqref{eq:rm-conditions}.
\end{theorem}

Сходимость требует, чтобы стратегия в пределе стала жадной, поэтому исследование должно уменьшаться со временем.  На практике часто используют постоянные $\varepsilon$ или $\alpha$, при которых SARSA сходится к окрестности $Q^*$, а не к $Q^*$ точно.

\subsection{Разобранный пример: хождение по скалам}

\begin{example}[Хождение по скалам]
\label{ex:cliff-walking}
\index{Cliff Walking}
Рассмотрим клеточный мир $4 \times 12$ с обрывом вдоль нижнего края (клетки $(3,1)$--$(3,10)$).  Агент начинает в $(3,0)$ и должен достичь цели в $(3,11)$.  Наступание на обрыв даёт вознаграждение $-100$ и возвращает агента на старт.  Все остальные переходы дают вознаграждение $-1$.  Эпизоды завершаются в цели или после 200 шагов.

SARSA с $\varepsilon=0.1$ обучается \emph{безопасному пути} вдоль верха сетки, держась подальше от обрыва.  Поскольку исследовательские действия иногда подталкивают агента к обрыву, SARSA оценивает клетки обрыва как опасные и избегает их.  Q-learning (раздел~\ref{sec:q-learning}), будучи алгоритмом на другой стратегии, оценивает каждое состояние при \emph{жадной} стратегии и поэтому находит кратчайший путь вдоль края обрыва\,---\,однако этот путь уязвим на практике, поскольку $\varepsilon$-жадное исследование иногда приводит агента к обрыву во время обучения.  Это канонический пример компромисса между текущей и другой стратегией: SARSA безопаснее во время обучения; Q-learning находит теоретически оптимальную стратегию.
\end{example}

% ===========================================================
\section{Q-learning: управление TD на другой стратегии}
\label{sec:q-learning}

\subsection{Обновление Q-learning}

Q-learning (Уоткинс, 1989)\index{Q-learning} непосредственно аппроксимирует оптимальную функцию ценности действия $Q^*$, используя уравнение оптимальности Беллмана в качестве цели.  Поскольку $Q^*(s,a) = \E[R_{t+1}+\gamma\max_{a'}Q^*(S_{t+1},a') \mid S_t=s,A_t=a]$, одновыборочное приближение даёт:

\begin{equation}
\label{eq:q-learning}
\boxed{Q(S_t,A_t) \leftarrow Q(S_t,A_t) + \alpha\Bigl(R_{t+1} + \gamma\max_{a'} Q(S_{t+1},a') - Q(S_t,A_t)\Bigr).}
\end{equation}

Ключевое отличие от SARSA~\eqref{eq:sarsa} --- оператор \textbf{max}.  Независимо от того, как фактически выбирается $A_{t+1}$ (по $\varepsilon$-жадной или любой другой поведенческой стратегии), цель всегда использует жадное действие в $S_{t+1}$.  Это делает Q-learning \textbf{алгоритмом на другой стратегии}\index{off-policy}: поведенческая стратегия собирает данные, а обновление оценивает жадную (целевую) стратегию.

\begin{algorithm}[ht]
\caption{Q-learning (управление TD на другой стратегии)}
\label{alg:q-learning}
\KwIn{Шаг $\alpha$, дисконт $\gamma$, исследование $\varepsilon$}
Инициализировать $Q(s,a)$ для всех $s,a$; $Q(s_{\mathrm{term}},\cdot)\leftarrow 0$\;
\For{каждый эпизод}{
  Инициализировать $S$\;
  \For{каждый шаг пока $S$ не терминальное}{
    Выбрать $A$ из $Q$ с помощью $\varepsilon$-жадной стратегии\;
    Выполнить действие $A$; наблюдать $R, S'$\;
    $Q(S,A) \leftarrow Q(S,A) + \alpha\bigl(R + \gamma\max_{a'}Q(S',a') - Q(S,A)\bigr)$\;
    $S \leftarrow S'$\;
  }
}
\end{algorithm}

\subsection{Сходимость Q-learning}

\begin{theorem}[Сходимость Q-learning, Уоткинс и Дайан, 1992]
\label{thm:q-learning-convergence}
\index{Q-learning!convergence}
В конечном МПР, если все пары состояние-действие посещаются бесконечно часто, вознаграждения ограничены и шаги удовлетворяют условиям Роббинса--Монро~\eqref{eq:rm-conditions}, то Q-learning сходится к $Q^*$ почти наверное.
\end{theorem}

Этот результат замечателен: он справедлив независимо от поведенческой стратегии, используемой для генерации переходов, при условии, что каждая пара $(s,a)$ посещается достаточно часто.  Q-learning является полностью внестратегийным методом с гарантиями сходимости в табличном случае.

\subsection{Смещение максимизации}
\label{subsec:max-bias}
\index{maximisation bias}

Тонкая, но важная проблема возникает из-за использования $\max_{a'} Q(S_{t+1},a')$ в качестве цели.  Когда $Q$ --- зашумлённая оценка, максимум зашумлённых значений является \emph{положительно смещённым} оценщиком истинного максимального значения.  Это связано с тем, что по неравенству Йенсена, применённому к выпуклой функции максимума:
\[
\E\!\left[\max_{a'}Q(S_{t+1},a')\right] \ge \max_{a'}\E[Q(S_{t+1},a')].
\]
Смещение может быть существенным, когда много действий имеют схожие истинные значения, но зашумлённые оценки.

\begin{example}[Смещение максимизации]
Предположим, что состояние $s'$ имеет три действия с истинными значениями $Q^*(s',a_1)=Q^*(s',a_2)=Q^*(s',a_3)=0$.  Оценки $\hat Q(s',a_i)\sim\cN(0,\sigma^2)$ --- независимый гауссовский шум.  Тогда $\E[\max_i\hat Q(s',a_i)] > 0$ --- положительное смещение, которое растёт с $\sigma^2$ и числом действий.  Q-learning использует это смещённое максимальное значение как цель, вызывая систематическое завышение оценок.
\end{example}

% ===========================================================
\section{Expected SARSA}
\label{sec:expected-sarsa}

\index{Expected SARSA}
SARSA использует фактическое следующее действие $A_{t+1}$ в своей цели, вводя дисперсию из-за случайности стратегии.  Q-learning устраняет эту дисперсию, беря жадное действие, но несёт смещение максимизации.  \textbf{Expected SARSA} (ван Сейен и др., 2009) выбирает средний путь: заменяет $Q(S_{t+1},A_{t+1})$ \emph{математическим ожиданием} по стратегии:

\begin{equation}
\label{eq:expected-sarsa}
Q(S_t,A_t) \leftarrow Q(S_t,A_t) + \alpha\Biggl(R_{t+1} + \gamma\sum_{a'}\pi(a'\mid S_{t+1})\,Q(S_{t+1},a') - Q(S_t,A_t)\Biggr).
\end{equation}

\begin{algorithm}[ht]
\caption{Expected SARSA}
\label{alg:expected-sarsa}
\KwIn{Шаг $\alpha$, дисконт $\gamma$, стратегия $\pi$ (или $\varepsilon$-жадная)}
Инициализировать $Q(s,a)$ для всех $s,a$\;
\For{каждый эпизод}{
  Инициализировать $S$\;
  \For{каждый шаг пока $S$ не терминальное}{
    Выбрать $A$ из $\pi(\cdot\mid S)$\;
    Выполнить действие $A$; наблюдать $R,S'$\;
    Вычислить $\bar Q(S') \leftarrow \sum_{a'}\pi(a'\mid S')\,Q(S',a')$\;
    $Q(S,A) \leftarrow Q(S,A) + \alpha\bigl(R + \gamma\,\bar Q(S') - Q(S,A)\bigr)$\;
    $S \leftarrow S'$\;
  }
}
\end{algorithm}

Expected SARSA устраняет дисперсию от случайного выбора $A_{t+1}$, поэтому его обновления менее зашумлены, чем у SARSA.  Его также можно использовать вне текущей стратегии: если $\pi$ является жадной стратегией (вероятность 1 на $\argmax_a Q(s,a)$, вероятность 0 иначе), Expected SARSA сводится в точности к Q-learning, поскольку
\[
\sum_{a'}\pi_{\mbox{жадная}}(a'\mid s')Q(s',a') = \max_{a'}Q(s',a').
\]
Таким образом, Expected SARSA обобщает как SARSA (на текущей стратегии, $\pi$ --- поведенческая), так и Q-learning (на другой стратегии, $\pi$ --- жадная).

% ===========================================================
\section{Double Q-learning}
\label{sec:double-q}

\index{Double Q-learning}
Double Q-learning (ван Хасселт, 2010) устраняет смещение максимизации, поддерживая \emph{две независимые} оценки функции ценности действия, $Q_1$ и $Q_2$.  Идея --- \emph{разделить} выбор действия и оценку действия:
\begin{itemize}[leftmargin=*]
  \item Использовать $Q_1$ для выбора жадного действия: $a^* = \argmax_{a'} Q_1(S_{t+1},a')$.
  \item Использовать $Q_2$ для оценки этого действия: цель $= R_{t+1} + \gamma Q_2(S_{t+1},a^*)$.
\end{itemize}
Поскольку $Q_1$ и $Q_2$ --- независимые оценки, использование одной для выбора и другой для оценки устраняет положительное смещение.  Важно, что роли $Q_1$ и $Q_2$ меняются с вероятностью $0.5$ на каждом шаге, чтобы обе оценки оставались актуальными.

\begin{algorithm}[ht]
\caption{Double Q-learning}
\label{alg:double-q}
\KwIn{Шаг $\alpha$, дисконт $\gamma$, исследование $\varepsilon$}
Инициализировать $Q_1(s,a)$ и $Q_2(s,a)$ для всех $s,a$\;
\For{каждый эпизод}{
  Инициализировать $S$\;
  \For{каждый шаг пока $S$ не терминальное}{
    $A \leftarrow \varepsilon$-жадная по $Q_1(S,\cdot)+Q_2(S,\cdot)$\;
    Наблюдать $R,S'$\;
    \eIf{с вероятностью $0.5$}{
      $a^* \leftarrow \argmax_{a'}Q_1(S',a')$\;
      $Q_1(S,A) \leftarrow Q_1(S,A) + \alpha\bigl(R+\gamma Q_2(S',a^*)-Q_1(S,A)\bigr)$\;
    }{
      $a^* \leftarrow \argmax_{a'}Q_2(S',a')$\;
      $Q_2(S,A) \leftarrow Q_2(S,A) + \alpha\bigl(R+\gamma Q_1(S',a^*)-Q_2(S,A)\bigr)$\;
    }
    $S \leftarrow S'$\;
  }
}
\end{algorithm}

Double Q-learning доказуемо уменьшает смещение максимизации и работает лучше, чем Q-learning, в средах с большим числом действий.  Эта идея позднее была масштабирована до глубокого ОП в алгоритме Double DQN (ван Хасселт, Гез и Сильвер, 2016), который использует онлайн-сеть для выбора действия и целевую сеть для оценки.

% ===========================================================
\section{Сравнение: SARSA, Q-learning и Expected SARSA}
\label{sec:td-control-comparison}

В таблице~\ref{tab:td-control-comparison} обобщены три основных алгоритма управления TD.

\begin{table}[ht]
\centering
\caption{Сравнение алгоритмов управления TD на текущей и другой стратегии.}
\label{tab:td-control-comparison}
\small
\begin{tabular}{@{}lccc@{}}
\toprule
\textbf{Свойство} & \textbf{SARSA} & \textbf{Q-learning} & \textbf{Expected SARSA}\\
\midrule
Цель & $R\!+\!\gamma Q(S',A')$ & $R\!+\!\gamma\max_{a'}Q(S',a')$ & $R\!+\!\gamma\!\sum_{a'}\pi(a'|S')Q(S',a')$\\
Стратегия & На текущей & На другой & Обе\\
Задача & $Q^\pi$ & $Q^*$ & $Q^\pi$ или $Q^*$\\
Дисперсия & Средняя & Низкая & Ниже, чем у SARSA\\
Смещение & Несмещённое & Смещение макс. & Несмещённое\\
Безопасно? & Да & Нет (обрыв) & Да\\
\bottomrule
\end{tabular}
\end{table}

\begin{figure}[ht]
\centering
\begin{tikzpicture}[x=0.65cm, y=0.65cm, >=stealth]
  % Нарисовать сетку
  \foreach \x in {0,...,11} {
    \foreach \y in {0,...,3} {
      \draw[gray!40] (\x,\y) rectangle (\x+1,\y+1);
    }
  }
  % Клетки обрыва (нижняя строка, столбцы 1-10)
  \foreach \x in {1,...,10} {
    \fill[red!60] (\x,0) rectangle (\x+1,1);
    \node[font=\tiny,white] at (\x+0.5,0.5) {О};
  }
  % Старт
  \fill[blue!60] (0,0) rectangle (1,1);
  \node[font=\tiny,white] at (0.5,0.5) {С};
  % Цель
  \fill[green!60!black] (11,0) rectangle (12,1);
  \node[font=\tiny,white] at (11.5,0.5) {Ц};
  % Граница сетки
  \draw[thick] (0,0) rectangle (12,4);
  % Путь SARSA (вдоль верха)
  \draw[->, very thick, blue!70!black, dashed]
    (0.5,0.5) -- (0.5,3.5) -- (11.5,3.5) -- (11.5,0.5);
  \node[blue!70!black,font=\small,above] at (6,3.8) {SARSA (безопасный путь)};
  % Путь Q-learning (вдоль края обрыва)
  \draw[->, very thick, orange!80!black]
    (0.5,0.5) -- (0.5,1.5) -- (11.5,1.5) -- (11.5,0.5);
  \node[orange!80!black,font=\small,below] at (6,-0.3) {Q-learning (оптимальный путь)};
  % Подписи
  \node[left,font=\small] at (0,2) {Сетка $4\times 12$};
\end{tikzpicture}
\caption{Клеточный мир «хождение по скалам».  Красные клетки --- обрыв (вознаграждение $-100$, сброс на старт).  SARSA (синий пунктир) обучается безопасному пути вдоль верха; Q-learning (сплошной оранжевый) обучается теоретически более короткому пути вдоль края обрыва, но этот путь уязвим при $\varepsilon$-жадном исследовании.}
\label{fig:cliff-walking}
\end{figure}

Пример «хождение по скалам» (пример~\ref{ex:cliff-walking} и рисунок~\ref{fig:cliff-walking}) иллюстрирует практическое следствие оценивания на текущей и другой стратегии.  Во время обучения с $\varepsilon$-жадным исследованием SARSA учитывает возможность случайных действий, толкающих агента к обрыву, и поэтому предпочитает более безопасный маршрут.  Q-learning оценивает состояния при жадной стратегии\,---\,которая никогда не сходит с края обрыва\,---\,поэтому обучается маршруту вдоль обрыва, оптимальному, но уязвимому.  При тестировании ($\varepsilon=0$) оба находят эквивалентные пути, но SARSA достигает более высокого среднего вознаграждения во время обучения.

% ===========================================================
\section{$n$-шаговые методы TD}
\label{sec:nstep-td}

\subsection{$n$-шаговая отдача}

Оба метода --- TD(0) и Монте-Карло --- являются частными случаями более общего семейства, параметризованного \emph{горизонтом просмотра вперёд} $n\in\{1,2,\dots\}$.  \textbf{$n$-шаговая отдача}\index{n-step return} определяется как

\begin{equation}
\label{eq:nstep-return}
G_t^{(n)} = R_{t+1} + \gamma R_{t+2} + \cdots + \gamma^{n-1}R_{t+n} + \gamma^n V(S_{t+n}),
\end{equation}
а соответствующее обновление:
\begin{equation}
\label{eq:nstep-update}
V(S_t) \leftarrow V(S_t) + \alpha\bigl(G_t^{(n)} - V(S_t)\bigr).
\end{equation}

\begin{proposition}
$n$-шаговая отдача удовлетворяет следующим свойствам:
\begin{enumerate}
  \item При $n=1$: $G_t^{(1)} = R_{t+1}+\gamma V(S_{t+1})$ --- цель TD(0).
  \item В эпизодической задаче с терминальным временем $T$: $G_t^{(T-t)} = G_t$ --- полная отдача Монте-Карло (при $V(S_T)=0$).
\end{enumerate}
\end{proposition}

\begin{proof}
(1) Подстановка $n=1$ в~\eqref{eq:nstep-return} даёт $G_t^{(1)} = R_{t+1}+\gamma V(S_{t+1})$.  (2) При $n=T-t$ и $V(S_T)=0$: $G_t^{(T-t)} = \sum_{k=0}^{T-t-1}\gamma^k R_{t+k+1} = G_t$.
\end{proof}

\subsection{Компромисс смещение--дисперсия в $n$-шаговом TD}

По мере увеличения $n$ $n$-шаговая отдача интерполирует между смещённой, но с низкой дисперсией целью TD(0) и несмещённой, но с высокой дисперсией отдачей МК:
\begin{itemize}[leftmargin=*]
  \item Малое $n$: больше бутстрэппинга, выше смещение, ниже дисперсия.
  \item Большое $n$: меньше бутстрэппинга, ниже смещение, выше дисперсия.
\end{itemize}
Оптимальное $n$ зависит от задачи: в средах с плотными вознаграждениями с низкой дисперсией малое $n$ работает хорошо; в средах с редкими вознаграждениями часто полезнее большее $n$.

\begin{figure}[ht]
\centering
\resizebox{\textwidth}{!}{%
\begin{tikzpicture}[x=1.1cm, y=1cm, >=stealth,
    box/.style={align=center, font=\footnotesize, text width=3.0cm}]
  \draw[->, thick] (0,0) -- (13,0) node[right] {$n$};
  \foreach \x/\lbl in {1/{$n=1$}, 6.5/{$1<n<T$}, 12/{$n=T-t$}}{
    \draw[thick] (\x,0.15) -- (\x,-0.15) node[below=6pt,font=\small]{\lbl};
  }
  \node[above=8pt,font=\bfseries] at (1,0) {TD(0)};
  \node[above=8pt,font=\bfseries] at (6.5,0) {$n$-шаговое TD};
  \node[above=8pt,font=\bfseries] at (12,0) {Монте-Карло};
  \node[box, below=1.2cm] at (1,0)   {Макс. бутстрэппинг\\Выс. смещение\\Низ. дисперсия};
  \node[box, below=1.2cm] at (6.5,0) {Частич. бутстрэп.\\Настр. горизонт\\Баланс bias--var};
  \node[box, below=1.2cm] at (12,0)  {Без бутстрэп.\\Низ. смещение\\Выс. дисперсия};
  \draw[very thin, gray, dashed] (1,-0.2) -- (1,-1.15);
  \draw[very thin, gray, dashed] (6.5,-0.2) -- (6.5,-1.15);
  \draw[very thin, gray, dashed] (12,-0.2) -- (12,-1.15);
  \draw[->,gray!60] (1,-3.5) -- (12,-3.5)
    node[midway,below=4pt,font=\scriptsize,align=center]{рост $n$: меньше бутстрэппинга, больше реального сигнала};
\end{tikzpicture}
}%
\caption{Спектр $n$-шаговых TD-методов.  TD(0) ($n=1$) и Монте-Карло ($n=T-t$) --- крайние случаи.  Промежуточные $n$ обеспечивают компромисс смещение--дисперсия, настраиваемый под задачу.}
\label{fig:nstep-spectrum}
\end{figure}

\subsection{$n$-шаговый SARSA}

Идея $n$-шагового подхода естественно распространяется на управление.  Цель $n$-шагового SARSA:
\[
G_t^{(n)} = R_{t+1} + \gamma R_{t+2} + \cdots + \gamma^{n-1}R_{t+n} + \gamma^n Q(S_{t+n},A_{t+n}),
\]
а обновление: $Q(S_t,A_t) \leftarrow Q(S_t,A_t)+\alpha(G_t^{(n)}-Q(S_t,A_t))$.  Этот алгоритм обладает гарантиями сходимости, аналогичными стандартному SARSA при подходящих условиях.

% ===========================================================
\section{TD($\lambda$) и следы приемлемости}
\label{sec:td-lambda}

Методы предыдущего раздела требуют фиксации одного горизонта $n$ заранее.  TD($\lambda$) вместо этого объединяет \emph{все} $n$-шаговые отдачи одновременно, взвешивая их геометрически так, что объединённая цель вычисляется онлайн с помощью одного дополнительного вектора --- \emph{следа приемлемости}.

\subsection{$\lambda$-отдача}

\begin{definition}[$\lambda$-отдача]
\label{def:lambda-return}
\index{lambda-return@$\lambda$-return}
Для $\lambda\in[0,1]$ \textbf{$\lambda$-отдача} --- это геометрически взвешенное среднее всех $n$-шаговых отдач:
\begin{equation}
\label{eq:lambda-return}
G_t^\lambda = (1-\lambda)\sum_{n=1}^{\infty}\lambda^{n-1}G_t^{(n)}.
\end{equation}
\end{definition}

Вес, присваиваемый $n$-шаговой отдаче, равен $(1-\lambda)\lambda^{n-1}$.  Эти веса положительны, геометрически убывают с ростом $n$ и в сумме равны единице:
\[
\sum_{n=1}^{\infty}(1-\lambda)\lambda^{n-1} = (1-\lambda)\cdot\frac{1}{1-\lambda} = 1.
\]
Таким образом, $G_t^\lambda$ является настоящей выпуклой комбинацией $n$-шаговых отдач.

\paragraph{Частные случаи.}
\begin{itemize}[leftmargin=*]
  \item $\lambda = 0$: $G_t^0 = G_t^{(1)} = R_{t+1}+\gamma V(S_{t+1})$ --- цель TD(0).  Весь вес на одношаговой отдаче.
  \item $\lambda = 1$: $G_t^1 = G_t$ --- полная отдача Монте-Карло.  Геометрические веса равномерны (в пределе дают полную отдачу).
\end{itemize}
Таким образом, TD($\lambda$) поглощает как TD(0), так и Монте-Карло как крайние случаи, как и $n$-шаговое TD\,---\,но с элегантной дифференцируемой интерполяцией.

\begin{figure}[ht]
\centering
\begin{tikzpicture}[x=0.7cm, y=2.5cm, >=stealth]
  \def\lam{0.6}
  \node[font=\small\bfseries] at (6.5,2.2) {Веса $\lambda$-отдачи $(1-\lambda)\lambda^{n-1}$ при $\lambda=0.6$};
  \foreach \n in {1,...,12}{
    \pgfmathsetmacro\w{(1-\lam)*(\lam^(\n-1))}
    \pgfmathsetmacro\wscaled{\w*5}
    \fill[blue!50!white] (\n-0.4,0) rectangle (\n+0.4,\wscaled);
    \draw[blue!70!black] (\n-0.4,0) rectangle (\n+0.4,\wscaled);
  }
  \draw[->,thick] (0.5,0) -- (13.5,0) node[right,font=\small] {$n$};
  \draw[->,thick] (0.5,0) -- (0.5,2.1) node[above,font=\small] {вес};
  \foreach \n in {1,2,3,4,5,6,7,8,9,10,11,12}{
    \draw ({\n},0.02) -- ({\n},-0.02);
    \node[below,font=\tiny] at (\n,-0.02) {$\n$};
  }
  \foreach \y/\lbl in {0.5/0.10, 1.0/0.20, 1.5/0.30, 2.0/0.40}{
    \draw (0.4,\y) -- (0.6,\y);
    \node[left,font=\tiny] at (0.4,\y) {\lbl};
  }
\end{tikzpicture}
\caption{Веса $(1-\lambda)\lambda^{n-1}$, присваиваемые каждой $n$-шаговой отдаче в $\lambda$-отдаче ($\lambda=0.6$).  Краткосрочные отдачи получают наибольший вес; веса геометрически убывают, поэтому влияние далёких отдач быстро уменьшается.}
\label{fig:lambda-weights}
\end{figure}

\subsection{Взгляд вперёд на TD($\lambda$)}
\label{subsec:forward-view}
\index{TD(lambda)!forward view}

\textbf{Взгляд вперёд} на TD($\lambda$) определяет цель обновления для состояния $S_t$ как $G_t^\lambda$ и применяет
\[
V(S_t) \leftarrow V(S_t) + \alpha\bigl(G_t^\lambda - V(S_t)\bigr).
\]
Это концептуально прозрачно: цель является конкретной, хорошо определённой величиной.  Однако вычисление $G_t^\lambda$ требует всей будущей траектории с момента $t$, поэтому взгляд вперёд требует ждать завершения эпизода до применения любого обновления.  Он поэтому не является онлайновым в операциональном смысле и эквивалентен однопроходному апостериорному обновлению по эпизоду.

Несмотря на непрактичность для онлайн-обучения, взгляд вперёд является \emph{золотым стандартом} для понимания того, что вычисляет TD($\lambda$).  Взгляд назад, который мы развиваем далее, математически эквивалентен взгляду вперёд, но работает онлайн.

\subsection{Следы приемлемости}
\label{subsec:eligibility-traces}
\index{eligibility trace}

Ключ к превращению TD($\lambda$) в онлайновый алгоритм --- \textbf{след приемлемости}\index{eligibility trace} $e_t(s)$, вектор по состояниям, который регистрирует, насколько «приемлемо» каждое состояние для кредитного обновления в момент $t$.

\begin{definition}[След приемлемости (накапливающий)]
\label{def:elig-trace}
\textbf{Накапливающий след приемлемости} для состояния $s$ инициализируется значением $e_0(s)=0$ и обновляется каждый шаг по правилу:
\begin{equation}
\label{eq:elig-trace}
e_t(s) = \gamma\lambda\, e_{t-1}(s) + \ind\{S_t = s\}.
\end{equation}
\end{definition}

След $e_t(s)$ увеличивается на 1 при каждом посещении состояния $s$ и убывает с коэффициентом $\gamma\lambda$ на каждом шаге времени.  Интуитивно он измеряет давность и частоту посещений $s$: состояние, посещаемое часто и недавно, имеет высокий след.

Альтернатива --- \textbf{заменяющий след приемлемости}\index{eligibility trace!replacing}:
\begin{equation}
\label{eq:elig-trace-replacing}
e_t(s) = \begin{cases} \gamma\lambda\, e_{t-1}(s) + 1 & \mbox{если } s = S_t,\\ \gamma\lambda\, e_{t-1}(s) & \mbox{иначе,}\end{cases}
\quad\mbox{ограниченный значением 1.}
\end{equation}
Заменяющие следы не позволяют следу для текущего состояния превышать 1, что может стабилизировать обучение в некоторых условиях.

\subsection{Взгляд назад на TD($\lambda$): онлайновый алгоритм}
\label{subsec:backward-view}
\index{TD(lambda)!backward view}

Во \textbf{взгляде назад} на каждом шаге вычисляется одна ошибка TD $\delta_t$, и \emph{все} состояния обновляются одновременно, взвешенные своими следами приемлемости:

\begin{equation}
\label{eq:td-lambda-update}
\boxed{V(s) \leftarrow V(s) + \alpha\,\delta_t\,e_t(s), \quad \forall s\in\cS.}
\end{equation}

Ошибка TD $\delta_t = R_{t+1}+\gamma V(S_{t+1})-V(S_t)$ транслируется во все состояния, но только состояния с ненулевым следом получают значимое обновление.  Недавно посещённые состояния\,---\,с высокими следами\,---\,получают большие обновления; состояния, не посещавшиеся недавно, имеют почти нулевые следы и практически не затрагиваются.  Это и есть механизм «присвоения кредита» следов приемлемости: хорошие или плохие исходы засчитываются пропорционально давности и частоте прошлых посещений.

\begin{algorithm}[ht]
\caption{TD($\lambda$) со следами приемлемости}
\label{alg:td-lambda}
\KwIn{Шаг $\alpha$, дисконт $\gamma$, параметр следа $\lambda$}
Инициализировать $V(s)\leftarrow 0$ для всех $s$\;
\For{каждый эпизод}{
  $e(s) \leftarrow 0$ для всех $s$\;
  Инициализировать $S$\;
  \For{каждый шаг $t$ пока $S$ не терминальное}{
    Выбрать $A\sim\pi(\cdot\mid S)$; наблюдать $R, S'$\;
    $\delta \leftarrow R + \gamma V(S') - V(S)$\;
    $e(S) \leftarrow e(S) + 1$ \quad(накапливающий след)\;
    \For{всех $s\in\cS$}{
      $V(s) \leftarrow V(s) + \alpha\,\delta\,e(s)$\;
      $e(s) \leftarrow \gamma\lambda\,e(s)$\;
    }
    $S \leftarrow S'$\;
  }
}
\end{algorithm}

На практике обновлять нужно только состояния с ненулевыми следами, поэтому внутренний цикл по всем $s$ обычно заменяется поддержкой множества активных следов.  Алгоритм требует $O(|\cS|)$ памяти на эпизод\,---\,одно значение следа на состояние.

\subsection{Эквивалентность взглядов вперёд и назад}

Следующая теорема является краеугольным камнем теории следов приемлемости.

\begin{theorem}[Эквивалентность взглядов вперёд и назад для TD($\lambda$)]
\label{thm:fb-equivalence}
\index{TD(lambda)!forward-backward equivalence}
Для автономной (обновляемой в конце эпизода) версии TD($\lambda$) с постоянным шагом $\alpha$ сумма всех обновлений $V(s)$ за эпизод одинакова при взгляде вперёд (верхний индекс fw) и при взгляде назад (верхний индекс bw).  То есть, определяя $\Delta_t^{\mathrm{fw}} V(s) = \alpha(G_t^\lambda - V(S_t))\ind\{S_t=s\}$ (вперёд) и $\Delta_t^{\mathrm{bw}}V(s) = \alpha\delta_t e_t(s)$ (назад):
\begin{equation}
\label{eq:fb-equiv}
\sum_{t=0}^{T-1}\Delta_t^{\mathrm{bw}}V(s) = \sum_{t=0}^{T-1}\Delta_t^{\mathrm{fw}}V(s), \quad \forall s\in\cS.
\end{equation}
\end{theorem}

\begin{proof}[Набросок доказательства]
Ключевое тождество:
\[
G_t^\lambda - V(S_t) = \sum_{k=t}^{T-1}(\gamma\lambda)^{k-t}\delta_k,
\]
которое выражает ошибку $\lambda$-отдачи как дисконтированную сумму будущих ошибок TD.  Для его проверки следует раскрыть $G_t^\lambda$ по определению~\eqref{eq:lambda-return}, записать каждую $G_t^{(n)}$ как телескопическую сумму ошибок TD (используя рекурсивно $G_t^{(n)} = \delta_t + \gamma G_{t+1}^{(n-1)}$), затем собрать слагаемые.  Полный алгебраический вывод можно найти у Саттона и Барто (2018, глава 12).

Установив это тождество, суммирование $\alpha\delta_t e_t(s)$ по $t$ и перестановка порядка суммирования даёт $\sum_t \alpha(G_t^\lambda-V(S_t))\ind\{S_t=s\}$, что завершает доказательство.
\end{proof}

Теорема~\ref{thm:fb-equivalence} подтверждает, что онлайновый инкрементный взгляд назад вычисляет в точности те же обновления, что и автономный взгляд вперёд.  Это глубокий результат: он означает, что нам не нужно выбирать между концептуальной ясностью (вперёд) и вычислительной эффективностью (назад).

\subsection{Практические соображения}

\begin{remark}[Затухание следа и эффективный горизонт]
Произведение $\gamma\lambda$ управляет эффективной скоростью затухания следов приемлемости.  При $\gamma\lambda = 0.9$ состояние, посещённое $k$ шагов назад, имеет след $0.9^k$, поэтому состояния, посещавшиеся более $\sim 20$ шагов назад, имеют пренебрежимо малый след.  На практике $\lambda$ часто настраивают совместно с $\gamma$ для установки желаемого эффективного горизонта.
\end{remark}

\begin{remark}[Вычислительная стоимость]
Наивная реализация алгоритма~\ref{alg:td-lambda} обновляет все $|\cS|$ следов на каждом шаге, что составляет $O(|\cS|)$ за шаг.  С табличными представлениями это осуществимо для умеренных пространств состояний.  При аппроксимации функций (глава~\ref{ch:value-approx}) след становится вектором той же размерности, что и вектор параметров, и стоимость за шаг составляет $O(d)$, где $d$ --- число параметров.
\end{remark}

\begin{remark}[SARSA($\lambda$) и Q($\lambda$)]
Следы приемлемости естественно распространяются на управление.  \textbf{SARSA($\lambda$)}\index{SARSA(lambda)@SARSA($\lambda$)} поддерживает след $e_t(s,a)$ по парам состояние-действие и использует ошибку TD SARSA $\delta_t = R_{t+1}+\gamma Q(S_{t+1},A_{t+1})-Q(S_t,A_t)$ для обновления $Q(s,a) \leftarrow Q(s,a)+\alpha\delta_t e_t(s,a)$.  Аналогично, \textbf{Q($\lambda$)}\index{Q(lambda)@Q($\lambda$)} применяет ошибку TD Q-learning со следами.  SARSA($\lambda$) особенно популярен на практике, поскольку сохраняет характер SARSA на текущей стратегии, включая многошаговое присвоение кредита следов приемлемости.
\end{remark}

\subsection{Разобранный пример: случайное блуждание с TD($\lambda$)}

\begin{example}[Влияние $\lambda$ на скорость обучения]
Вернёмся к пятисостоянию случайного блуждания из примера~\ref{ex:random-walk}.  Запустив TD($\lambda$) с различными значениями $\lambda$ для фиксированного числа эпизодов и измерив СКО относительно истинных значений, наблюдаем:
\begin{itemize}[leftmargin=*]
  \item $\lambda=0$ (TD(0)): медленное обучение; сигнал вознаграждения требует многих эпизодов для распространения в далёкие состояния.
  \item $\lambda=0.5$: следы обеспечивают распространение вознаграждения из правого терминала назад к состояниям, посещённым в течение последних нескольких шагов в том же эпизоде, существенно ускоряя обучение.
  \item $\lambda=0.9$: быстрое начальное обучение, но слегка более зашумлённое, поскольку следы придают значительный вес далёким состояниям, которые могли не вносить вклада в итоговое вознаграждение.
  \item $\lambda=1$ (МК): наибольшая дисперсия за эпизод, но при хорошо настроенном шаге суммарное обучение по эпизодам может быть конкурентоспособным.
\end{itemize}
В экспериментах Саттона и Барто промежуточные значения $\lambda\approx 0.7$--$0.9$ часто достигают наименьшего СКО для этой задачи.  Оптимальное $\lambda$ зависит от шага и структуры задачи.
\end{example}

% ===========================================================
\section{Обобщённая оценка преимущества (GAE)}
\label{sec:gae}

\index{Generalised Advantage Estimation (GAE)}
\index{advantage function}
Обобщённая оценка преимущества (Шульман и др., 2016) --- современный метод вычисления оценок преимущества с низкой дисперсией для использования в методах градиента стратегии (глава~\ref{ch:pg}).  Он построен непосредственно на механизме TD($\lambda$) и представляет собой убедительное применение теории, разработанной в этой главе.

\subsection{Функция преимущества}

\begin{definition}[Функция преимущества]
\label{def:advantage}
\index{advantage function}
\textbf{Функция преимущества} стратегии $\pi$ равна:
\begin{equation}
\label{eq:advantage}
A^\pi(s,a) = Q^\pi(s,a) - V^\pi(s).
\end{equation}
\end{definition}

Преимущество $A^\pi(s,a)$ измеряет, насколько действие $a$ лучше среднего действия при стратегии $\pi$ в состоянии $s$.  По определению:
\begin{itemize}[leftmargin=*]
  \item $A^\pi(s,a) > 0$: действие $a$ лучше среднего; стратегия должна увеличить его вероятность.
  \item $A^\pi(s,a) < 0$: действие $a$ хуже среднего; уменьшить его вероятность.
  \item $A^\pi(s,a) = 0$: действие $a$ работает ровно на уровне среднего.
\end{itemize}
Заметим, что $\E_{a\sim\pi(\cdot|s)}[A^\pi(s,a)] = 0$ для всех $s$, то есть преимущество центрировано.

\paragraph{Почему преимущество важно для градиентов стратегии.}
Теорема о градиенте стратегии (глава~\ref{ch:pg}) утверждает, что градиент ожидаемой отдачи по параметрам стратегии $\theta$ равен
\[
\nabla_\theta J(\theta) = \E_\pi\bigl[\nabla_\theta\log\pi_\theta(A_t\mid S_t)\cdot\Psi_t\bigr],
\]
где $\Psi_t$ может быть отдачей $G_t$, функцией ценности действия $Q^\pi(S_t,A_t)$ или\,---\,что особенно важно\,---\,преимуществом $A^\pi(S_t,A_t)$.  Использование преимущества вместо сырой отдачи или функции ценности действия резко снижает дисперсию: базовая линия $V^\pi(S_t)$ вычитается, устраняя всю вариацию, обусловленную состоянием (на которое стратегия не может влиять).

\subsection{$n$-шаговые оценки преимущества}

На практике $V^\pi$ и $A^\pi$ неизвестны.  При наличии обученного критика $V_\phi \approx V^\pi$ естественным оценщиком одношагового преимущества является ошибка TD:
\begin{equation}
\label{eq:td-advantage}
\hat A_t^{(1)} = \delta_t = R_{t+1} + \gamma V_\phi(S_{t+1}) - V_\phi(S_t).
\end{equation}
В более общем виде $n$-шаговая оценка преимущества:
\begin{equation}
\label{eq:nstep-advantage}
\hat A_t^{(n)} = \sum_{l=0}^{n-1}\gamma^l R_{t+l+1} + \gamma^n V_\phi(S_{t+n}) - V_\phi(S_t).
\end{equation}
Заметим, что $\hat A_t^{(n)} = G_t^{(n)} - V_\phi(S_t)$, где $G_t^{(n)}$ --- $n$-шаговая отдача~\eqref{eq:nstep-return}.

\subsection{GAE: преимущество с весами по $\lambda$}

\begin{definition}[Обобщённая оценка преимущества]
\label{def:gae}
\index{GAE}
\textbf{Обобщённая оценка преимущества} $\hat A_t^{\mathrm{GAE}(\gamma,\lambda)}$ --- это экспоненциально взвешенное среднее $n$-шаговых оценок преимущества:
\begin{equation}
\label{eq:gae}
\hat A_t^{\mathrm{GAE}(\gamma,\lambda)} = \sum_{l=0}^{\infty}(\gamma\lambda)^l\,\delta_{t+l},
\end{equation}
где $\delta_{t+l} = R_{t+l+1} + \gamma V_\phi(S_{t+l+1}) - V_\phi(S_{t+l})$ --- ошибка TD в момент $t+l$.
\end{definition}

\begin{proposition}
GAE является $\lambda$-взвешенным средним $n$-шаговых оценок преимущества:
\begin{equation}
\label{eq:gae-lambda-avg}
\hat A_t^{\mathrm{GAE}(\gamma,\lambda)} = (1-\lambda)\sum_{n=1}^{\infty}\lambda^{n-1}\hat A_t^{(n)}.
\end{equation}
\end{proposition}

\begin{proof}
Используем ключевое тождество $\hat A_t^{(n)} = \sum_{l=0}^{n-1}(\gamma)^l\delta_{t+l}$ (телескопическое), которое следует из
\[
\hat A_t^{(n)} = \sum_{l=0}^{n-1}\gamma^l R_{t+l+1} + \gamma^n V_\phi(S_{t+n}) - V_\phi(S_t)
= \sum_{l=0}^{n-1}\gamma^l\delta_{t+l}.
\]
Тогда
\begin{align*}
(1-\lambda)\sum_{n=1}^{\infty}\lambda^{n-1}\hat A_t^{(n)}
&= (1-\lambda)\sum_{n=1}^{\infty}\lambda^{n-1}\sum_{l=0}^{n-1}\gamma^l\delta_{t+l}\\
&= (1-\lambda)\sum_{l=0}^{\infty}\gamma^l\delta_{t+l}\sum_{n=l+1}^{\infty}\lambda^{n-1}\\
&= (1-\lambda)\sum_{l=0}^{\infty}\gamma^l\delta_{t+l}\cdot\frac{\lambda^l}{1-\lambda}\\
&= \sum_{l=0}^{\infty}(\gamma\lambda)^l\delta_{t+l}.
\end{align*}
\end{proof}

\subsection{Частные случаи и компромисс смещение--дисперсия}

\begin{itemize}[leftmargin=*]
  \item $\lambda = 0$: $\hat A_t^{\mathrm{GAE}(\gamma,0)} = \delta_t = R_{t+1}+\gamma V_\phi(S_{t+1})-V_\phi(S_t)$ --- одношаговое преимущество TD с минимальной дисперсией и максимальным смещением (полная зависимость от $V_\phi$).
  \item $\lambda = 1$: $\hat A_t^{\mathrm{GAE}(\gamma,1)} = \sum_{l=0}^{\infty}\gamma^l\delta_{t+l} = G_t - V_\phi(S_t)$ --- оценка преимущества Монте-Карло с минимальным смещением и максимальной дисперсией.
  \item Промежуточные $\lambda$: экспоненциально убывающие веса усредняют краткосрочные (низкая дисперсия, высокое смещение) и долгосрочные (высокая дисперсия, низкое смещение) оценки.
\end{itemize}

\subsection{Связь с TD($\lambda$)}

Структура GAE~\eqref{eq:gae} идентична механизму $\lambda$-отдачи.  Фактически, если рассматривать $V_\phi(S_t)$ как базовую линию, то
\[
\hat A_t^{\mathrm{GAE}(\gamma,\lambda)} = G_t^\lambda - V_\phi(S_t),
\]
где $G_t^\lambda$ --- $\lambda$-отдача, вычисленная с критиком $V_\phi$.  Таким образом, GAE --- это просто цель TD($\lambda$) минус текущая оценка ценности.  Взгляд на следы приемлемости предоставляет эффективный способ вычисления GAE онлайн: поддерживать след $e_t$ и накапливать слагаемые $e_t\,\delta_t$.

На практике (например, в алгоритмах PPO и TRPO) GAE вычисляется обратным проходом по собранному роллауту длины $T$:
\begin{align*}
A_T &= 0,\\
A_t &= \delta_t + \gamma\lambda\,A_{t+1}, \quad t = T-1, T-2, \ldots, 0,
\end{align*}
что является линейным рекуррентным соотношением, вычисляемым за время $O(T)$.

\begin{example}[Вычисление GAE для короткого роллаута]
Предположим $\gamma=0.99$, $\lambda=0.95$, и роллаут из 4 шагов даёт ошибки TD $\delta_0=0.5,\;\delta_1=-0.3,\;\delta_2=0.8,\;\delta_3=0.1$ (при $A_4=0$).  Обратное вычисление даёт:
\begin{align*}
A_3 &= 0.1 + 0 = 0.1,\\
A_2 &= 0.8 + 0.99\times 0.95\times 0.1 = 0.8 + 0.094 = 0.894,\\
A_1 &= -0.3 + 0.99\times 0.95\times 0.894 = -0.3 + 0.840 = 0.540,\\
A_0 &= 0.5 + 0.99\times 0.95\times 0.540 = 0.5 + 0.508 = 1.008.
\end{align*}
Оценка преимущества $A_0 = 1.008$ для первого шага учитывает не только $\delta_0 = 0.5$, но и нижестоящий положительный сигнал $\delta_2 = 0.8$, соответствующим образом дисконтированный на $({\gamma\lambda})^2$.
\end{example}

% ===========================================================
\section{Анализ смещения, дисперсии и сходимости}
\label{sec:td-bias-variance}

\subsection{Формальный анализ смещения}

Цель TD(0) $R_{t+1}+\gamma V_t(S_{t+1})$ смещена, когда $V_t \ne V^\pi$.  Разложим смещение следующим образом.  Пусть $V_t$ обозначает оценку ценности на итерации~$t$, и пусть $b_t(s) = V_t(s) - V^\pi(s)$ --- ошибка приближения на итерации~$t$.  Тогда
\begin{align*}
\E_\pi[R_{t+1}+\gamma V_t(S_{t+1})\mid S_t=s]
&= V^\pi(s) + \gamma\,\E_\pi[b_t(S_{t+1})\mid S_t=s].
\end{align*}
Таким образом, смещение цели TD равно $\gamma$, умноженному на ожидаемую ошибку приближения в следующем состоянии.  По величине это строго меньше ошибки приближения в $s$ (масштабированной на $\gamma < 1$), что и является механизмом, благодаря которому TD-обновления притягиваются к $V^\pi$.  По мере того как $V_t\to V^\pi$, смещение исчезает, что согласуется с теоремой~\ref{thm:td-error-zero-mean}.

Монте-Карло, напротив, использует несмещённую цель $G_t$, поэтому его смещение равно нулю независимо от $V_t$.

\subsection{Анализ дисперсии}

Дисперсия цели TD(0) $R_{t+1}+\gamma V(S_{t+1})$ ограничена:
\[
\Var[R_{t+1}+\gamma V(S_{t+1})\mid S_t=s] \le \sigma_R^2 + \gamma^2\,\Var[V(S_{t+1})\mid S_t=s],
\]
где $\sigma_R^2$ --- дисперсия вознаграждения.  Когда $V$ --- гладкая функция (например, нейронная сеть), $\Var[V(S_{t+1})]$ мало.  Для сравнения, отдача МК $G_t = \sum_{k=0}^{\infty}\gamma^k R_{t+k+1}$ накапливает дисперсию с каждого шага.  Поскольку вознаграждения на разных шагах, вообще говоря, коррелированы через общую траекторию, полное выражение включает члены с ковариациями:
\begin{align*}
\Var[G_t\mid S_t=s] &= \sum_{k=0}^{\infty}\gamma^{2k}\Var[R_{t+k+1}\mid S_t=s] \\
&\quad + 2\!\!\sum_{0\le j<k}\!\!\gamma^{j+k}\Cov[R_{t+j+1},R_{t+k+1}\mid S_t=s] \\
&\gg \Var[R_{t+1}+\gamma V(S_{t+1})].
\end{align*}
Этот разрыв в дисперсии растёт с длиной эпизода и коэффициентом дисконтирования, делая методы TD существенно более выборочно эффективными, чем МК, во многих практических условиях.

\subsection{Разложение среднеквадратичной ошибки}

Для любого оценщика $\hat y$ цели $y^* = \E[Y]$ среднеквадратичная ошибка раскладывается как:
\begin{equation}
\label{eq:bias-variance-decomp}
\mbox{МКО}(\hat y) = \underbrace{\bigl(\E[\hat y] - y^*\bigr)^2}_{\mbox{Смещение}^2} + \underbrace{\Var[\hat y]}_{\mbox{Дисперсия}}.
\end{equation}
Это разложение применяется к цели TD (смещённой, с низкой дисперсией) и отдаче МК (несмещённой, с высокой дисперсией).  Минимизация МКО требует баланса обоих слагаемых, что именно и делает TD($\lambda$) настройкой $\lambda$.

\subsection{Сравнительная таблица}

\begin{table}[ht]
\centering
\caption{Свойства смещения, дисперсии и сходимости методов предсказания.}
\label{tab:bias-variance-comparison}
\begin{tabular}{lcccc}
\toprule
\textbf{Метод} & \textbf{Смещение} & \textbf{Дисперсия} & \textbf{Онлайн?} & \textbf{Сходимость}\\
\midrule
Монте-Карло & Нулевое & Высокая (растёт с $T$) & Нет & Да (к $V^\pi$)\\
TD(0) & Высокое нач., $\to 0$ & Низкая & Да & Да (Теор.~\ref{thm:td0-convergence})\\
$n$-шаговое TD & Среднее & Средняя & Полуонлайн & Да\\
TD($\lambda$) & Настраиваемое ($\lambda$) & Настраиваемая ($\lambda$) & Да & Да\\
\bottomrule
\end{tabular}
\end{table}

\subsection{Смертельная триада}
\index{deadly triad}

Табличные TD-методы гарантированно сходятся при мягких условиях.  Однако сочетание трёх ингредиентов ведёт к потенциальной расходимости:
\begin{enumerate}
  \item \textbf{Аппроксимация функций} (например, линейные или нейросетевые представления $V$).
  \item \textbf{Бутстрэппинг} (TD-обновления с использованием текущей оценки в качестве цели).
  \item \textbf{Обучение на другой стратегии} (обучение по данным, сгенерированным другой стратегией).
\end{enumerate}
Это сочетание\,---\,известное как \emph{смертельная триада}\,---\,может привести к расходимости оценок ценности до бесконечности даже в простых линейных примерах (контрпример Бэрда, 1995).  Монте-Карло избегает триады за счёт отказа от бутстрэппинга; методы на текущей стратегии избегают её, оставаясь на текущей стратегии; некоторые алгоритмы (например, методы градиентного TD, emphatic TD) решают её посредством тщательного алгоритмического проектирования.  Эти вопросы подробно рассматриваются в главе~\ref{ch:value-approx}.

% ===========================================================
\section{Итоги}
\label{sec:td-summary}

Обучение с временной разностью объединяет традиции Монте-Карло и динамического программирования.  Из единственной идеи бутстрэппинга\,---\,использования текущих оценок ценности в качестве целей\,---\,возникает целое семейство эффективных, онлайновых, безмодельных алгоритмов.  В таблице~\ref{tab:td-summary} приводится полный обзор всех методов, введённых в этой главе.

\begin{table}[ht]
\centering
\caption{Сводка методов временной разности, введённых в этой главе.}
\label{tab:td-summary}
\footnotesize
\begin{tabular}{@{}lllccc@{}}
\toprule
\textbf{Алгоритм} & \textbf{Тип} & \textbf{Стр-гия} & \textbf{Цель} & \textbf{Онлайн} & \textbf{Кредит}\\
\midrule
TD(0) & Предск. & Фикс. $\pi$ & $R+\gamma V(S')$ & Да & 1\\
$n$-шаг. TD & Предск. & Фикс. $\pi$ & $G^{(n)}$ & Полу & $n$\\
TD($\lambda$) & Предск. & Фикс. $\pi$ & $G^\lambda$ & Да & Следы\\
SARSA & Упр-е & Текущ. & $R+\gamma Q(S',A')$ & Да & 1\\
SARSA($\lambda$) & Упр-е & Текущ. & $G^\lambda_Q$ & Да & Следы\\
Q-learning & Упр-е & Другой & $R+\gamma\max_{a'}Q(S',a')$ & Да & 1\\
Expected SARSA & Упр-е & Обе & $R+\gamma\bar Q(S')$ & Да & 1\\
Double Q-learning & Упр-е & Другой & Раздел. макс. & Да & 1\\
\bottomrule
\end{tabular}
\end{table}

Ключевые выводы этой главы:
\begin{itemize}[leftmargin=*]
  \item Ошибка TD $\delta_t = R_{t+1}+\gamma V(S_{t+1})-V(S_t)$ --- фундаментальный кредитный сигнал в обучении с подкреплением; при сходимости её среднее равно нулю (теорема~\ref{thm:td-error-zero-mean}).
  \item Методы TD обменивают смещение на дисперсию: бутстрэппинг вводит смещение, но снижает дисперсию по сравнению с Монте-Карло.
  \item SARSA работает на текущей стратегии и безопасен при исследовании; Q-learning работает на другой стратегии и напрямую обучает оптимальной стратегии.
  \item TD($\lambda$) со следами приемлемости элегантно заполняет разрыв между одношаговым TD и Монте-Карло, при этом взгляды вперёд и назад доказуемо эквивалентны.
  \item GAE применяет механизм $\lambda$-отдачи к оценке преимущества в методах градиента стратегии, обеспечивая чёткий компромисс смещение--дисперсия, параметризованный $\lambda$.
  \item Сходимость в табличном случае хорошо изучена; сочетание аппроксимации функций, бутстрэппинга и данных с другой стратегии (смертельная триада) может вызвать расходимость.
\end{itemize}


% ===========================================================
\section*{Упражнения}
\addcontentsline{toc}{section}{Упражнения}

\begin{exercise}
\label{ex:td-error-basic}
Предположим $V(S_t)=3$, $R_{t+1}=1$, $\gamma=0.9$, $V(S_{t+1})=4$.
\begin{enumerate}[label=(\alph*)]
  \item Вычислите ошибку TD $\delta_t$.
  \item Примените одно обновление TD(0) с шагом $\alpha=0.1$.  Каково новое значение $V(S_t)$?
  \item Интерпретируйте знак $\delta_t$: было ли состояние переоценено или недооценено?
\end{enumerate}
\end{exercise}

\begin{exercise}
\label{ex:td-error-zero-mean}
Пусть $V = V^\pi$ --- точная функция ценности.  Дайте подробное доказательство того, что $\E_\pi[\delta_t \mid S_t=s]=0$ для всех $s$, явно указав, на каком шаге используется уравнение Беллмана для $V^\pi$.  Затем объясните, почему это свойство \emph{не} выполняется при $V \ne V^\pi$.
\end{exercise}

\begin{exercise}
\label{ex:sarsa-vs-qlearn}
Рассмотрите среду «хождение по скалам» из примера~\ref{ex:cliff-walking}.
\begin{enumerate}[label=(\alph*)]
  \item Объясните точными словами, почему SARSA с $\varepsilon=0.1$ предпочитает путь вдали от обрыва, тогда как Q-learning --- путь вдоль края.
  \item Какой алгоритм найдёт лучшую политику при $\varepsilon \to 0$?  Докажите свой ответ, апеллируя к теоремам о сходимости.
\end{enumerate}
\end{exercise}

\begin{exercise}
\label{ex:double-q-bias}
Рассмотрите состояние $s'$ с $n$ действиями, для которых истинные значения $Q^*(s',a_i)=0$ для всех $i$, а оценки $\hat Q(s',a_i)\sim\cN(0,\sigma^2)$ независимы.
\begin{enumerate}[label=(\alph*)]
  \item Покажите, что $\E[\max_i\hat Q(s',a_i)] > 0$, и вычислите этот предел для $n=2$ аналитически.
  \item Объясните, почему Double Q-learning устраняет это смещение в ожидании.
\end{enumerate}
\end{exercise}

\begin{exercise}
\label{ex:td-lambda}
Для TD($\lambda$) с накапливающими следами:
\begin{enumerate}[label=(\alph*)]
  \item Запишите условие на $\lambda$ и $\gamma$, при котором след состояния $s$, посещённого $k$ шагов назад, составляет менее $0.01$ от его первоначального значения.
  \item Покажите, что при $\lambda=1$ и завершении эпизода обновления TD($\lambda$) в конце эпизода совпадают с обновлениями Монте-Карло.
\end{enumerate}
\end{exercise}

\begin{exercise}
\label{ex:gae}
Для GAE с параметрами $\gamma=0.99$ и $\lambda=0.95$ на роллауте из 5 шагов с TD-ошибками $\delta_t = (0.2, -0.1, 0.5, 0.3, -0.4)$ (в момент $t=0,1,2,3,4$ соответственно):
\begin{enumerate}[label=(\alph*)]
  \item Вычислите оценку GAE $\hat A_0$ с помощью обратного рекуррентного соотношения.
  \item Каков предел $\hat A_0$ при $\lambda \to 0$?  При $\lambda \to 1$?
\end{enumerate}
\end{exercise}
